\section{Introduction}

Narrative similarity is often treated as a semantic matching problem, where models prioritize lexical overlap, topical closeness, or stylistic resemblance. However, many downstream tasks require sensitivity to \emph{what happens} in a story and \emph{how events unfold} over time. Two narratives may use different wording and details yet still share the same underlying event progression, while semantically similar texts may diverge in causal and temporal structure.

This distinction matters for practical systems. In retrieval-augmented generation (RAG), a retriever that captures narrative structure can return context that is behaviorally relevant rather than merely topically related. In recommendation systems, users often care about plot dynamics (e.g., setup, escalation, reversal, resolution), not only keywords or genres. As large language models continue to improve, direct LLM-based scoring is increasingly effective but still expensive to scale across large corpora. A strong embedding model that approximates structure-aware judgments offers a lower-latency, lower-cost alternative for large-scale indexing, retrieval, and ranking.

In this work, we focus on learning story embeddings that are explicitly sensitive to \emph{structural similarity}. Our goal is not simply to improve generic semantic similarity, which is already well served by strong off-the-shelf embedding models. Instead, we ask whether an embedding model can be trained to place stories close together when they share event progression, temporal organization, and narrative dynamics, even when they differ in topic, entities, wording, or style.

A central challenge is that structural similarity is not a standard supervised target. To ground the task, we run a user study that asks how people perceive structural similarity between narratives. This study is important because it determines what our training signal should emphasize: not just whether two stories discuss related content, but whether their event sequences align in a way that humans recognize as structurally similar. We use these judgments to calibrate a structural similarity model that combines event alignment and aligned-event semantics.

The resulting calibrated model acts as a teacher for embedding training. We apply it to larger ASQ-derived story-pair data to obtain pseudo-labels, then fine-tune neural embedding models so cosine similarity reflects the teacher's structure-aware scores. We evaluate the resulting embeddings across several complementary datasets and task formats, including pair separation, retrieval, ranking, synthetic anchor comparisons, and external narrative-similarity benchmarks. This lets us test whether structural supervision transfers beyond the pairwise training setup into retrieval and ranking settings that better resemble practical applications.
